\documentclass[11pt]{article}
\usepackage[margin=1in]{geometry}
\usepackage[T1]{fontenc}
\usepackage{lmodern}
\usepackage{microtype}
\usepackage{amsmath,amssymb,amsthm}
\usepackage{booktabs}
\usepackage{xcolor}
\usepackage[hidelinks]{hyperref}
\usepackage{xurl}
\usepackage{enumitem}
\setlist{nosep, leftmargin=*}
\emergencystretch=3em
\newtheorem{proposition}{Proposition}

\title{Gated DeltaNet on AWS Inferentia2 Without a Triangular Solve:\\
Exact Static Graphs for Hybrid Qwen3.5/Qwen3.8 Models}
\author{Saksham Adhikari}
\date{Technical report, November 2026 \\ \small Not peer reviewed.}

\begin{document}
\maketitle

\begin{abstract}
Qwen3.5 and Qwen3.8 interleave three Gated DeltaNet (GDN) linear-attention layers with every full-attention layer. The
chunked GDN algorithm inverts a unit lower-triangular $64\times64$ matrix in every chunk, and the PyTorch reference does this
with \texttt{torch.linalg.solve\_triangular}. Existing Neuron ports replace that step with NKI kernels that run only on
Trainium2, and Neuron SDK 2.29 removed NKI-kernel support for Trainium1 and Inferentia2 in NxD Inference. We first show what happens to
\texttt{solve\_triangular} on Inferentia2 under SDK 2.28, the last SDK for that chip. It \emph{compiles without error and is
wrong}. torch-xla has no device lowering for the operator, so tracing runs it on the host and freezes its output into the
compiled graph as a constant. The result is exact on the tracing input and has median relative error 0.56 on any other input.
Traced inside the whole model, it compiles and then crashes at run time. We then remove the operator. The inverse is computed
exactly by $\log_2 64=6$ block-doubling steps of dense matrix products under constant masks. Padded positions use a gate and
step size that provably leave the recurrent state unchanged. All state is explicit and fixed-shape, and no tensor exceeds
rank~4. On an Inferentia2 chip the doubling inverse is accurate to $1.6\times10^{-7}$ on real systems it was not traced with,
and it is $2.9\times$ faster than the only other formulation that generalises, the row-by-row loop used as a fallback in
Hugging Face Transformers. The whole Qwen3.5-0.8B model compiles to one prefill and one decode graph. With bf16 weights it
reproduces 32 of 32 greedy tokens of the fp32 reference on both test prompts, with KL divergence at most $7.3\times10^{-4}$,
including on prompts it was not traced with. Prefill of a $2\times256$ batch takes 41--52\,ms and a batch-2 decode step 21.7\,ms.
All device runs used Spot instances and cost about \$0.20 in total.
\end{abstract}

\section{Introduction}
Hybrid decoders replace most softmax-attention layers with linear-attention or state-space layers. These keep a fixed-size
recurrent state instead of a growing KV cache~\cite{lieber2024jamba,minimax2025,nemotronh2025,kimi2025linear}. The Qwen family
adopted Gated DeltaNet~\cite{yang2025gdn} in Qwen3-Next~\cite{qwen3next}, then in Qwen3.5~\cite{qwen35} and
Qwen3.8~\cite{qwen38card,qiu2026qwen38next}. Qwen3.8-27B has 64 layers in a 3:1 GDN-to-attention pattern.

Prefill uses the chunkwise-parallel form of the delta rule~\cite{yang2024deltanet}. Each chunk of $C=64$ tokens needs the
inverse of a unit lower-triangular matrix (the UT transform~\cite{joffrain2006ut}). GPU libraries implement it in Triton
kernels~\cite{fla}. The Hugging Face reference calls \texttt{torch.linalg.solve\_triangular}, with a 63-step loop of in-place
row updates as a fallback~\cite{hftransformers}. The Neuron documentation reports unsupported operators with error
NCC\_EVRF001~\cite{neuronevrf001}. The usual explanation for the difficulty is that the accelerator ``needs fixed tensor
sizes'', but that is not the problem here: the chunk is $64\times64$ in every call.

The Neuron ports of this family that we found replace the inverse with NKI kernels on Trainium2: NxD Inference's Qwen3.5
contrib models~\cite{nxdicontrib}, the NKI library's \texttt{gdn\_cte}~\cite{nkilibgdn} and a community
kernel~\cite{jburtoft}. NxD Inference's notes also record that the plain PyTorch chunk path, which builds rank-5 tensors,
crashes neuronx-cc code generation (NCC\_INLA001)~\cite{nxdicontrib}. Since SDK 2.29, NxD Inference no longer supports NKI
kernels on Trainium1 or Inferentia2, and users there are told to stay on SDK 2.28~\cite{neuron229}. On those chips, the inverse has to be expressed with operators
the compiler handles on its own.

\paragraph{Contributions.}
\begin{enumerate}
\item \textbf{A silent failure mode.} On Inferentia2 with SDK 2.28 and \texttt{torch\_neuronx.trace},
\texttt{solve\_triangular} is not rejected. It is evaluated on the host at trace time and its output is baked into the NEFF.
The graph is exact on the tracing input and wrong on every other input (median relative error 0.56). A whole-model graph
containing it compiles and then segfaults. A test that evaluates on its tracing input cannot detect this; we evaluate every
graph on inputs it was not traced with (Section~\ref{sec:solve}).
\item \textbf{An exact matmul-only inverse} with a proof. It costs 12 dense $C\times C$ products per chunk. On device it is
accurate to fp32 rounding on unseen real systems, and $2.9\times$ faster to run and $5\times$ faster to compile than the
generalising row loop (Sections~\ref{sec:inverse},~\ref{sec:device}).
\item \textbf{One static graph per phase} for the whole hybrid model. It uses a padding rule ($g=0,\ \beta=0$) proven to leave
the state unchanged, explicit fixed-shape state, rank-$\le4$ tensors and matmul prefix sums (Section~\ref{sec:graph}).
\item \textbf{Measurements on a real checkpoint, CPU and Inferentia2.} These include the first Gated DeltaNet hybrid we know
of to run on Inferentia2. They also show that random-weight tests miss the instability of the Neumann alternative
(Sections~\ref{sec:eval},~\ref{sec:device}).
\item \textbf{A reproducible, cost-bounded recipe.} One-time Spot instances upload their results and terminate themselves.
Code and raw results are released (Section~\ref{sec:repro}).
\end{enumerate}

\section{Background}
\paragraph{Gated delta rule.} For one head with keys and queries $k_t,q_t\in\mathbb{R}^{d_k}$ (L2-normalised), values
$v_t\in\mathbb{R}^{d_v}$, decay $\alpha_t=e^{g_t}\in(0,1]$ and step size $\beta_t\in(0,1)$, the state $S_t\in\mathbb{R}^{d_k\times d_v}$
evolves as
\begin{equation}
S_t = \alpha_t S_{t-1} + k_t\,\beta_t\big(v_t - \alpha_t S_{t-1}^{\top}k_t\big)^{\top}, \qquad o_t = S_t^{\top} q_t .
\label{eq:rec}
\end{equation}
Qwen3.5/3.8 compute $\beta_t=\sigma(\cdot)$ and $g_t=-e^{A}\,\mathrm{softplus}(\cdot+\delta)$ per value head from the hidden
state, after a depthwise causal convolution (kernel 4) and SiLU on $[q;k;v]$. A gated RMSNorm follows the core.

\paragraph{Chunked form.} Within a chunk of $C$ positions, let $\gamma_i=\sum_{j\le i} g_j$ and
$\Gamma_{ij}=e^{\gamma_i-\gamma_j}$ for $i\ge j$ (0 otherwise). The UT system is $T = I + L$ with
$L = \mathrm{tril}\big((\beta\odot K)K^{\top}\odot\Gamma,\,-1\big)$. Then
\begin{equation}
U = T^{-1}(\beta\odot V), \quad W = T^{-1}(\beta\odot e^{\gamma}\odot K), \quad
V' = U - W S, \quad O = (e^{\gamma}\odot Q)S + (QK^{\top}\odot\Gamma)\,V',
\label{eq:chunk}
\end{equation}
followed by $S \leftarrow e^{\gamma_C} S + (e^{\gamma_C-\gamma}\odot K)^{\top} V'$. Everything in~\eqref{eq:chunk} is a dense
product except $T^{-1}$.

\paragraph{Conditioning.} Under these conditions (unit-norm keys, $\beta,\alpha\in(0,1]$), every entry of $T^{-1}$ is bounded by
1 in magnitude~\cite{sobczyk2026tri}. Sobczyk et al.\ find power-series (Neumann) inversion unstable in fp32 at $C\ge64$, and
recursive inversion stable. Classical analysis of blocked triangular inversion is in~\cite{ducroz1992,higham2002}.

\section{What \texttt{solve\_triangular} does on Inferentia2}\label{sec:solve}
\texttt{torch\_neuronx.trace} records the model through torch-xla and hands the resulting HLO to neuronx-cc. We traced a module
that returns $(I+L)^{-1}$ through \texttt{solve\_triangular} and dumped the HLO that torch-xla produced for it. The HLO is the
identity on a single parameter: there is no triangular-solve, loop or arithmetic instruction. torch-xla had no lowering for
the operator, ran it on the host, and passed its \emph{result} to the device as data. During tracing, that data becomes a constant
of the compiled graph.

The compiler therefore reports success. Table~\ref{tab:kernel} shows the consequence: on the 256 systems used for tracing,
the output matches fp64 to $5.5\times10^{-7}$, but on 256 other real systems from the same checkpoint, the median relative error is
0.56 and the maximum 1.44. The graph returns the tracing input's inverse whatever it is given. Inside the whole model, where
the solve's input depends on the prompt, the same trace compiled (one graph split around the host computation) and then
crashed the process with a segmentation fault before producing output. The rejection reported elsewhere (NCC\_EVRF001) is
not what a torch-neuronx user sees; what they see is a graph that passes a same-input check.

\section{A matmul-only exact inverse}\label{sec:inverse}
Let $D_b$ be the block-diagonal part of $T$ with $b\times b$ diagonal blocks, so $D_1=I$ and $D_C=T$. Let $O_b = T\odot M_b$,
where the constant 0/1 mask $M_b$ selects the lower-left $b\times b$ sub-block inside each diagonal $2b\times2b$ block. Then
$D_{2b}=D_b+O_b$.

\begin{proposition}
If $X_b = D_b^{-1}$, then $X_{2b} = X_b - X_b O_b X_b$ equals $D_{2b}^{-1}$ exactly.
\end{proposition}
\begin{proof}
$D_{2b}(X_b - X_bO_bX_b) = I + O_bX_b - O_bX_b - O_bX_bO_bX_b = I - (O_bX_bO_b)X_b$, using $D_bX_b=I$. Within each $2b$ block,
$O_b$ is nonzero only in (lower-half rows, upper-half columns). $X_b$ is block-diagonal at size $b$, so $O_bX_b$ has the same
support. In $(O_bX_b)O_b$ the summation index must be an upper-half column of the first factor and a lower-half row of the
second at the same time, so the product is zero.
\end{proof}

Starting from $X_1=I$, six steps give $T^{-1}$ for $C=64$. Each step uses two dense $C\times C$ products and an elementwise
mask. There is no data-dependent indexing, no division and no loop over rows. Every intermediate $X_b$ is the exact inverse of
a principal block-diagonal part of $T$, so its entries obey the same unit bound~\cite{sobczyk2026tri}. By contrast, the
Neumann product $(I-L)(I+L^2)(I+L^4)\cdots$ forms powers $L^{2^j}$ whose entries grow with $\lVert L\rVert$ before they cancel,
which is what fails in fp32 (Table~\ref{tab:inverse}).

\paragraph{Cost.} The dense form spends $6\times 2\times 2C^3\approx 6.3$\,MFLOP per chunk and head, against about 11\,MFLOP
for the rest of the chunk with $d_k=d_v=128$. Only about 0.2\,MFLOP of the doubling work is non-zero. A kernel could skip the
zeros, but the dense form keeps the graph to operators that every backend lowers. Its HLO has 102 instructions of 11 kinds,
mostly \texttt{dot}, \texttt{multiply} and \texttt{subtract}. The row loop needs 2280 instructions with
\texttt{dynamic-update-slice} (Table~\ref{tab:kernel}).

\section{One static graph per phase}\label{sec:graph}
\paragraph{Padding that leaves the state unchanged.}
Prompts are right-padded to a bucket length that is a multiple of $C$. At padded positions we set $\beta_t=0$ and $g_t=0$
($\alpha_t=1$), and zero the convolution input.
\begin{proposition}
With $\beta_t=0$, $g_t=0$ at every padded position (all after the real tokens), the final state equals that of the unpadded
sequence, and every real-position output is unchanged.
\end{proposition}
\begin{proof}
In~\eqref{eq:rec}, $\beta_t=0,\alpha_t=1$ gives $S_t=S_{t-1}$. In the chunked form, a padded row $i$ of $L$ is
$\beta_i(\cdot)=0$, so row $i$ of $T^{-1}$ is $e_i^{\top}$. Row $i$ of $\beta\odot V$ and of $\beta\odot e^\gamma\odot K$ is
zero, so rows $i$ of $U$ and $W$, and hence of $V'$, are zero, and the padded row adds nothing to the state update. Because
$g_i=0$, $\gamma$ does not change after the last real token, so the chunk decay $e^{\gamma_C}$ and the key scaling of the real
rows are those of the unpadded chunk. Real outputs depend only on earlier positions by causality.
\end{proof}
The final convolution window must hold the last three \emph{real} inputs of each row. We gather it at index
$\ell+\{0,1,2\}$ from the input left-padded by three zeros, where $\ell$ is the row's length, so batches with mixed prompt
lengths are exact. Full-attention layers need no change: causal masking already hides later padded keys.

\paragraph{Explicit state.} Prefill maps $(\texttt{ids}[B,T],\ \ell[B])$ to the last-token logits and a flat tuple of state
tensors. Decode maps $(\texttt{tok}[B],\ \texttt{pos}[B],\ \text{state})$ to logits and the new state. The state holds a KV cache
$[B,h_{kv},N_{\max},256]$ per full-attention layer, a convolution window $[B,d_{conv},3]$ and an fp32 recurrent state
$[B,h_v,128,128]$ per GDN layer. Decode writes the KV slot with a one-hot \texttt{where} blend and masks attention by comparing
against \texttt{pos}; there is no scatter or dynamic slice.

\paragraph{Rank $\le4$, matmul prefix sums, no control flow.} Batch and heads are folded into one dimension inside the chunk
computation, so no tensor exceeds rank~4; rank~5 is what triggers NCC\_INLA001~\cite{nxdicontrib}. The within-chunk prefix sum
$\gamma$ is a product with a constant upper-triangular matrix of ones, so it runs on the Tensor engine rather than as a scan.
The loop over chunks is unrolled to $T/C$ steps at trace time. For Qwen3.8-27B the per-sequence GDN state is
$48\times48\times128\times128\times4$\,B $\approx 151$\,MB in fp32, and the KV cache adds 64\,KiB per token in bf16.

\section{Evaluation on CPU}\label{sec:eval}
We compare against Hugging Face Transformers 5.24 (\texttt{Qwen3\_5}, eager attention, its PyTorch GDN path with
\texttt{solve\_triangular}) in fp32 on CPU. We use the real Qwen3.5-0.8B checkpoint (24 layers, 18 GDN). Qwen3.8 has no
official checkpoint under 27B and shares \texttt{model\_type}~\texttt{qwen3\_5}. The batch holds prompts of 5 and 167 tokens,
right-padded to a 256-token bucket.

\begin{table}[t]
\centering\small
\begin{tabular}{rrrrr}
\toprule
GDN layer & $\max\lVert L\rVert_F$ & $\max|T^{-1}|$ & doubling (ours) & Neumann product \\
\midrule
0  & 6.6  & 1 & $1.7\times10^{-7}$ & $2.0\times10^{-5}$ \\
5  & 29.6 & 1 & $1.2\times10^{-7}$ & $9.8\times10^{4}$ \\
8  & 41.6 & 1 & $1.9\times10^{-7}$ & $1.3\times10^{10}$ \\
18 & 35.2 & 1 & $1.8\times10^{-7}$ & $6.7\times10^{8}$ \\
22 & 31.4 & 1 & $1.6\times10^{-7}$ & $8.6\times10^{9}$ \\
\midrule
all 18 layers & $\le 41.6$ & 1 & $\le 2.6\times10^{-7}$ & up to $1.3\times10^{10}$ \\
\bottomrule
\end{tabular}
\caption{Inverse of every real UT system in the prefill, CPU fp32: max entry error relative to an fp64
\texttt{solve\_triangular}, worst chunk and head. The measured $\max|T^{-1}|=1$ matches the unit bound
of~\cite{sobczyk2026tri}. The Neumann product is wrong wherever $\lVert L\rVert_F\gtrsim15$.}
\label{tab:inverse}
\end{table}

\begin{table}[t]
\centering\small
\begin{tabular}{p{0.34\linewidth}p{0.29\linewidth}p{0.29\linewidth}}
\toprule
Check & Qwen3.5-0.8B (real weights) & 27B head layout (random) \\
\midrule
Prefill last-token logits vs.\ HF (unpadded) & max $|\Delta|$ $3.2\times10^{-5}$, same argmax & max $|\Delta|$ $2.0\times10^{-5}$, same argmax \\
Greedy decode, 24 steps vs.\ HF \texttt{generate} & 24/24 identical; KL $\le3.7\times10^{-7}$ & 24/24 identical; KL $\le1.0\times10^{-6}$ \\
\texttt{torch.export} prefill / decode & 43 / 33 aten ops, max rank 4 & 43 / 33 aten ops, max rank 4 \\
Forbidden ops (solve, linalg, nonzero, \texttt{item}) & none & none \\
\bottomrule
\end{tabular}
\caption{CPU agreement with the reference. The second column uses a random-weight model with Qwen3.8-27B's head layout (48 GDN
value heads per 16 key heads, 24/4 attention heads, head dim 256, partial RoPE 0.25) to exercise the grouped-value path.}
\label{tab:e2e}
\end{table}

Tables~\ref{tab:inverse} and~\ref{tab:e2e} show the inverse is accurate to fp32 rounding on every real system and that the
whole model agrees with the reference. On the random-weight model, $\lVert L\rVert_F\le2$ and the Neumann product is also
accurate. A test suite built only on random weights would have accepted an inverse that fails on the real
checkpoint~\cite{jburtoft}.

\section{Evaluation on Inferentia2}\label{sec:device}
\paragraph{Setup.} One Inferentia2 chip (two NeuronCores, 32\,GB), on Spot \texttt{inf2.xlarge} and, when that type had no
Spot capacity, \texttt{inf2.8xlarge} (same chip, more host CPU). We used the AWS Neuron 2.28 DLAMI: torch-neuronx
2.9.0.2.12, neuronx-cc 2.23. Graphs were compiled with \texttt{torch\_neuronx.trace}~\cite{neurontrace}, one prefill graph per bucket~\cite{nxdibucketing}, and \texttt{--auto-cast=none}, so the fp32
parts of the recurrence stay fp32. Weights are bf16 and inlined into the NEFF. The reference is the same static model in fp32
on the host, which matches Transformers to $3\times10^{-5}$ (Table~\ref{tab:e2e}). Every graph is evaluated on its tracing
input and on inputs it was not traced with.

\begin{table}[t]
\centering\small
\begin{tabular}{lrrrrr}
\toprule
formulation & HLO instr. & compile & warm call & err.\ (traced input) & err.\ (unseen inputs) \\
\midrule
\texttt{solve\_triangular} & 7 (identity) & 4.3\,s & 0.53\,ms & $5.5\times10^{-7}$ & \textbf{1.44} (median 0.56) \\
row loop (HF fallback) & 2280 & 26.9\,s & 4.05\,ms & $4.1\times10^{-7}$ & $1.6\times10^{-7}$ \\
doubling (ours) & 102 & 5.3\,s & \textbf{1.41\,ms} & $2.3\times10^{-7}$ & $\mathbf{1.6\times10^{-7}}$ \\
\bottomrule
\end{tabular}
\caption{Inverse alone on Inferentia2: 256 real UT systems from Qwen3.5-0.8B, the hardest by $\lVert L\rVert_F$ (up to 41.6),
fp32. Errors are the max entry error relative to fp64, worst system. ``Unseen'' is the next 256 hardest systems, evaluated
with the graph traced on the first 256. The solve's warm time is a host-to-device copy of a constant.}
\label{tab:kernel}
\end{table}

\paragraph{Inverse.} Table~\ref{tab:kernel}. Of the formulations that compute on the device, the doubling inverse is the most
accurate and the fastest: $2.9\times$ faster per call and $5.1\times$ faster to compile than the row loop, with a NEFF
$6.8\times$ smaller (247\,KB vs.\ 1.7\,MB).

\begin{table}[t]
\centering\small
\begin{tabular}{lrr}
\toprule
& prompt A (5 tok.) & prompt B (167 tok.) \\
\midrule
KL(fp32 $\Vert$ device) at prefill & $7.3\times10^{-4}$ & $4.0\times10^{-4}$ \\
max / mean per-step KL, 32 steps (teacher-forced) & $7.3\times10^{-4}$ / $1.3\times10^{-4}$ & $4.0\times10^{-4}$ / $4.4\times10^{-5}$ \\
argmax agreement, teacher-forced & 32/32 & 32/32 \\
identical greedy prefix vs.\ fp32 & 32/32 & 32/32 \\
\midrule
& prompt C (10 tok.) & prompt D (90 tok.) \\
\midrule
prefill KL, prompts not traced with & $5.8\times10^{-4}$ & $2.4\times10^{-4}$ \\
argmax equal & yes & yes \\
\bottomrule
\end{tabular}
\caption{Whole model on Inferentia2 (bf16 weights, fp32 GDN state) vs.\ the fp32 host reference. The top block uses the batch the
graphs were traced with. The bottom block is a different batch run through the same prefill graph.}
\label{tab:device}
\end{table}

\paragraph{Whole model.} Both phases compile as single graphs: prefill $[2\times256]$ in 346\,s and decode $[2\times1]$ in
496\,s on \texttt{inf2.xlarge} (279\,s for prefill on \texttt{inf2.8xlarge}). Each NEFF is 1.6\,GB with the weights inlined.
Table~\ref{tab:device} gives the fidelity. Greedy decoding on the device reproduces all 32 tokens of the fp32 reference for both
prompts, and the per-step KL stays below $10^{-3}$, the size expected from bf16 weights. The prefill graph is equally accurate on
a batch it was not traced with, which is the check that \texttt{solve\_triangular} fails.

\paragraph{Latency.} Prefill of the $2\times256$ bucket (172 real tokens) takes 52\,ms when it also returns the state, and 41\,ms
for logits only. A batch-2 decode step takes 21.7\,ms (p90 21.9\,ms), or 92 tokens/s. That step time includes copying the whole
state, 51\,MB each way, between host and device, because the traced decode graph takes and returns the state as ordinary
tensors. Keeping the state resident on the device (input/output aliasing) is the obvious next step and would lower it. We
did not profile the share of device time spent in the inverse inside the full model.

\paragraph{Cost.} Three Spot runs, about 25, 11 and 23 minutes, at \$0.09/h (\texttt{inf2.xlarge}) and \$0.31--0.37/h
(\texttt{inf2.8xlarge}) in us-east-1, cost about \$0.20 in total. One run was reclaimed by Spot and repeated. Each instance
downloads the checkpoint, runs, uploads results to S3 and terminates itself, with a hard power-off as a backstop.

\section{Scaling to Qwen3.8-27B}
Qwen3.8-27B is 54\,GB in bf16 and needs tensor parallelism. TP=8 on \texttt{inf2.24xlarge} (four chips, 128\,GB of HBM in use)
divides the 16 GDN key heads (2 key and 6 value heads per core) and the 24 query heads, and replicates each of the 4 KV heads
twice. GDN heads are independent, so the recurrence needs no collectives beyond the usual all-reduce after the output
projections. We describe this configuration but did not run it. The chunk loop unrolls $T/C$ times in each of 48 GDN
layers, so compile time grows linearly with the bucket. Long contexts call for moderate buckets with chunked prefill across
calls, which the explicit state interface already supports, or an NKI kernel for the scan.

\section{Related work}
The delta rule's chunkwise form and its WY/UT representation come from parallel
DeltaNet~\cite{yang2024deltanet,bischof1987wy,joffrain2006ut}; Gated DeltaNet adds the decay gate~\cite{yang2025gdn}.
Related linear-attention and SSM layers include~\cite{katharopoulos2020,schlag2021,yang2024gla,dao2024mamba2}. For the
triangular inverse, FLA combines 16-row forward substitution with Schur merges~\cite{fla}. Sobczyk et al.\ analyse direct and
iterative inversion on NPUs~\cite{sobczyk2026tri}. A multiplication-only approximate inverse has been proposed for quantized GDN
on mobile NPUs~\cite{zhang2026goodenough}; ours is exact. On Neuron, the existing ports are NKI-based and target
Trainium2~\cite{nxdicontrib,nkilibgdn,jburtoft}. Serving systems manage hybrid state alongside paged KV
caches~\cite{vllmhybrid,sglanghybrid}.

\section{Limitations}
\begin{itemize}
\item Only Qwen3.5-0.8B ran on device, at one bucket (256) and batch 2. Qwen3.8-27B was validated only through its head layout
on random weights, on CPU.
\item The solve-triangular finding is for torch-neuronx tracing on SDK 2.28. Other front ends, such as NxD Inference's model
builder, or other SDK versions may behave differently.
\item The decode latency includes host round-trips of the state. Neither phase was profiled at the engine level.
\item The dense doubling inverse spends FLOPs on zero blocks, and a $64\times64$ product underuses the $128\times128$ Tensor
engine.
\item The multimodal encoder and the multi-token-prediction head are out of scope.
\end{itemize}

\section{Reproducibility}\label{sec:repro}
Code, launch scripts and raw JSON results are in the accompanying repository:
\begin{itemize}
\item \texttt{static\_qwen.py} (the model) and \texttt{test\_static.py} (CPU checks);
\item \texttt{neuron/run\_inf2.py} and \texttt{neuron/solve\_vs\_doubling.py} (device runs);
\item \texttt{infra/inf2\_spot.sh} (Spot launcher);
\item \texttt{results/} (device runs 20261123T183627Z, 20261123T190820Z and 20261123T194739Z).
\end{itemize}
CPU checks used PyTorch 2.15, Transformers 5.24 and Python 3.12.

\section{Conclusion}
On Inferentia2, \texttt{solve\_triangular} in Gated DeltaNet fails silently: it compiles, but its output is frozen at trace time.
Six masked doubling steps of dense products replace it exactly, run on the device, stay accurate to fp32 rounding on real
checkpoints where the Neumann product fails by ten orders of magnitude, and are $2.9\times$ faster than the row-loop
alternative. With state-neutral padding, explicit fixed-shape state and rank-4 tensors, the whole hybrid model becomes one static
graph per phase. On an Inferentia2 chip it reproduces the fp32 reference token for token. This route works on Inferentia2 and
Trainium1, where NKI-based ports are no longer supported.


\begin{thebibliography}{99}
\bibitem{bischof1987wy}
C.~Bischof and C.~Van~Loan.
\newblock The WY representation for products of Householder matrices.
\newblock \emph{SIAM Journal on Scientific and Statistical Computing}, 8(1):s2--s13, 1987.

\bibitem{dao2024mamba2}
T.~Dao and A.~Gu.
\newblock Transformers are SSMs: Generalized models and efficient algorithms through structured state space duality.
\newblock In \emph{Proceedings of ICML}, 2024. arXiv:2405.21060.

\bibitem{ducroz1992}
J.~J. Du~Croz and N.~J. Higham.
\newblock Stability of methods for matrix inversion.
\newblock \emph{IMA Journal of Numerical Analysis}, 12(1):1--19, 1992. (LAPACK Working Note 27.)

\bibitem{fla}
S.~Yang and Y.~Zhang.
\newblock FLA: A Triton-based library for hardware-efficient implementations of linear attention mechanism.
\newblock Software, \url{https://github.com/fla-org/flash-linear-attention} (\texttt{fla/ops/utils/solve\_tril.py}), accessed
November 2026. Not peer reviewed.

\bibitem{hftransformers}
Hugging Face.
\newblock \texttt{transformers}: \texttt{models/qwen3\_5/modeling\_qwen3\_5.py}, \texttt{torch\_chunk\_gated\_delta\_rule}.
\newblock Software, \url{https://github.com/huggingface/transformers}, main branch, accessed November 2026. Not peer reviewed.

\bibitem{higham2002}
N.~J. Higham.
\newblock \emph{Accuracy and Stability of Numerical Algorithms}.
\newblock SIAM, 2nd edition, 2002.

\bibitem{jburtoft}
jburtoft (Hugging Face user).
\newblock \texttt{qwen35-gdn-deltanet-neuron-kernels} (formerly \texttt{qwen35-deltanet-tkg-full}): NKI kernels for Qwen3.5 Gated
DeltaNet, v15.
\newblock Software, \url{https://huggingface.co/kernels/jburtoft/qwen35-deltanet-tkg-full}, accessed November 2026. Community
artifact; not peer reviewed.

\bibitem{joffrain2006ut}
T.~Joffrain, T.~M. Low, E.~S. Quintana-Ort\'i, R.~van~de~Geijn, and F.~G. Van~Zee.
\newblock Accumulating Householder transformations, revisited.
\newblock \emph{ACM Transactions on Mathematical Software}, 32(2):169--179, 2006.

\bibitem{katharopoulos2020}
A.~Katharopoulos, A.~Vyas, N.~Pappas, and F.~Fleuret.
\newblock Transformers are RNNs: Fast autoregressive transformers with linear attention.
\newblock In \emph{Proceedings of ICML}, 2020. arXiv:2006.16236.

\bibitem{kimi2025linear}
Kimi Team.
\newblock Kimi Linear: An expressive, efficient attention architecture.
\newblock arXiv:2510.26692, 2025.

\bibitem{lieber2024jamba}
O.~Lieber et~al.
\newblock Jamba: A hybrid Transformer-Mamba language model.
\newblock arXiv:2403.19887, 2024.

\bibitem{minimax2025}
MiniMax (A.~Li et~al.).
\newblock MiniMax-01: Scaling foundation models with lightning attention.
\newblock arXiv:2501.08313, 2025.

\bibitem{nemotronh2025}
NVIDIA (A.~Blakeman et~al.).
\newblock Nemotron-H: A family of accurate and efficient hybrid Mamba-Transformer models.
\newblock arXiv:2504.03624, 2025.

\bibitem{neuron229}
AWS Neuron.
\newblock Neuron 2.29.0 release notes (NxD Inference component) and Neuron 2.30 announcements.
\newblock \url{https://awsdocs-neuron.readthedocs-hosted.com/en/latest/release-notes/2.29.0.html}, accessed November 2026.

\bibitem{neuronevrf001}
AWS Neuron.
\newblock Neuron compiler error code \texttt{NCC\_EVRF001}: unsupported operator.
\newblock \url{https://awsdocs-neuron.readthedocs-hosted.com/en/v2.28.0/compiler/error-codes/EVRF001.html}, accessed November 2026.

\bibitem{neurontrace}
AWS Neuron.
\newblock PyTorch Neuron (\texttt{torch-neuronx}) tracing API for inference, SDK v2.28.0.
\newblock
\url{https://awsdocs-neuron.readthedocs-hosted.com/en/v2.28.0/frameworks/torch/torch-neuronx/api-reference-guide/inference/api-torch-neuronx-trace.html},
accessed November 2026.

\bibitem{nkilibgdn}
AWS Neuron.
\newblock NKI library, \texttt{nkilib.experimental.gdn} (\texttt{gdn\_cte}, \texttt{gdn\_cte\_schur}, \texttt{gdn\_cte\_neumann}) and
README.
\newblock Software, \url{https://github.com/aws-neuron/nki-library}, accessed November 2026. Not peer reviewed.

\bibitem{nxdibucketing}
AWS Neuron.
\newblock NxD Inference API reference guide: bucketing (\texttt{context\_encoding\_buckets}, \texttt{token\_generation\_buckets}).
\newblock \url{https://awsdocs-neuron.readthedocs-hosted.com/en/latest/libraries/nxd-inference/api-guides/api-guide.html}, accessed
November 2026.

\bibitem{nxdicontrib}
AWS Neuron and contributors.
\newblock NxD Inference contrib models \texttt{Qwen3.5-2B} and \texttt{Qwen3.5-35B-A3B}; pull request \#173 (Qwen3.6-27B, open).
\newblock Software, \url{https://github.com/aws-neuron/neuronx-distributed-inference}, accessed November 2026. Not peer reviewed.

\bibitem{qiu2026qwen38next}
Z.~Qiu, Z.~Wang, X.~Li, et~al.\ (36 authors).
\newblock On the design of Qwen3.8-Next architecture: Evaluation, efficiency, and training stability.
\newblock arXiv:2608.30320, 2026. (The text refers to the model as Qwen3.8-Flash-Next.)

\bibitem{qwen35}
Qwen Team.
\newblock Qwen3.5: Towards native multimodal agents.
\newblock Blog, \url{https://qwen.ai/blog?id=qwen3.5}, February 2026; model card \texttt{Qwen/Qwen3.5-397B-A17B}.

\bibitem{qwen38card}
Qwen Team.
\newblock \texttt{Qwen/Qwen3.8-27B} model card and \texttt{config.json}.
\newblock \url{https://huggingface.co/Qwen/Qwen3.8-27B}, August 2026, accessed November 2026.

\bibitem{qwen3next}
Qwen Team.
\newblock Qwen3-Next (blog) and \texttt{Qwen/Qwen3-Next-80B-A3B-Instruct} model card.
\newblock \url{https://huggingface.co/Qwen/Qwen3-Next-80B-A3B-Instruct}, September 2025, accessed November 2026.

\bibitem{schlag2021}
I.~Schlag, K.~Irie, and J.~Schmidhuber.
\newblock Linear Transformers are secretly fast weight programmers.
\newblock In \emph{Proceedings of ICML}, 2021. arXiv:2102.11174.

\bibitem{sglanghybrid}
SGLang contributors.
\newblock Hybrid model prefix-cache checkpoints (\texttt{--hybrid-cache-checkpoints-level}).
\newblock Documentation, \url{https://github.com/sgl-project/sglang} (\texttt{docs/docs/advanced\_features/hybrid\_cache\_checkpoints.mdx}),
accessed November 2026. Not peer reviewed.

\bibitem{sobczyk2026tri}
A.~Sobczyk, G.~Gottardo, C.~K. Matzoros, M.~De~Vita, F.~Skogh, A.~Zouzias, and J.~Zhuang.
\newblock Fast and stable triangular inversion for delta-rule linear transformers.
\newblock In \emph{Advances in Neural Information Processing Systems}, 2026. arXiv:2605.21325 (v2, extended version).

\bibitem{vllmhybrid}
vLLM contributors.
\newblock Hybrid KV cache manager (design document).
\newblock \url{https://github.com/vllm-project/vllm} (\texttt{docs/design/hybrid\_kv\_cache\_manager.md}), accessed November 2026. Not
peer reviewed.

\bibitem{yang2024deltanet}
S.~Yang, B.~Wang, Y.~Zhang, Y.~Shen, and Y.~Kim.
\newblock Parallelizing linear transformers with the delta rule over sequence length.
\newblock In \emph{Advances in Neural Information Processing Systems}, 2024. arXiv:2406.06484.

\bibitem{yang2024gla}
S.~Yang, B.~Wang, Y.~Shen, R.~Panda, and Y.~Kim.
\newblock Gated linear attention transformers with hardware-efficient training.
\newblock In \emph{Proceedings of ICML}, 2024. arXiv:2312.06635.

\bibitem{yang2025gdn}
S.~Yang, J.~Kautz, and A.~Hatamizadeh.
\newblock Gated delta networks: Improving Mamba2 with delta rule.
\newblock In \emph{International Conference on Learning Representations (ICLR)}, 2025. arXiv:2412.06464.

\bibitem{zhang2026goodenough}
L.~Zhang, Y.~Ren, K.~Zhang, T.~Liu, L.~Ge, D.~Li, M.~H. Langston, Y.~Huang, W.~W. Zeng, and L.~Zhang.
\newblock When good enough is optimal: Multiplication-only matrix inversion approximation for quantized Gated DeltaNet.
\newblock Poster, AdaptFM Workshop at ICML 2026. arXiv:2606.06034.

\end{thebibliography}
\end{document}
